\documentclass{letter}\usepackage[]{graphicx}\usepackage[]{xcolor}
% maxwidth is the original width if it is less than linewidth
% otherwise use linewidth (to make sure the graphics do not exceed the margin)
\makeatletter
\def\maxwidth{ %
  \ifdim\Gin@nat@width>\linewidth
    \linewidth
  \else
    \Gin@nat@width
  \fi
}
\makeatother

\definecolor{fgcolor}{rgb}{0.345, 0.345, 0.345}
\newcommand{\hlnum}[1]{\textcolor[rgb]{0.686,0.059,0.569}{#1}}%
\newcommand{\hlstr}[1]{\textcolor[rgb]{0.192,0.494,0.8}{#1}}%
\newcommand{\hlcom}[1]{\textcolor[rgb]{0.678,0.584,0.686}{\textit{#1}}}%
\newcommand{\hlopt}[1]{\textcolor[rgb]{0,0,0}{#1}}%
\newcommand{\hlstd}[1]{\textcolor[rgb]{0.345,0.345,0.345}{#1}}%
\newcommand{\hlkwa}[1]{\textcolor[rgb]{0.161,0.373,0.58}{\textbf{#1}}}%
\newcommand{\hlkwb}[1]{\textcolor[rgb]{0.69,0.353,0.396}{#1}}%
\newcommand{\hlkwc}[1]{\textcolor[rgb]{0.333,0.667,0.333}{#1}}%
\newcommand{\hlkwd}[1]{\textcolor[rgb]{0.737,0.353,0.396}{\textbf{#1}}}%
\let\hlipl\hlkwb

\usepackage{framed}
\makeatletter
\newenvironment{kframe}{%
 \def\at@end@of@kframe{}%
 \ifinner\ifhmode%
  \def\at@end@of@kframe{\end{minipage}}%
  \begin{minipage}{\columnwidth}%
 \fi\fi%
 \def\FrameCommand##1{\hskip\@totalleftmargin \hskip-\fboxsep
 \colorbox{shadecolor}{##1}\hskip-\fboxsep
     % There is no \\@totalrightmargin, so:
     \hskip-\linewidth \hskip-\@totalleftmargin \hskip\columnwidth}%
 \MakeFramed {\advance\hsize-\width
   \@totalleftmargin\z@ \linewidth\hsize
   \@setminipage}}%
 {\par\unskip\endMakeFramed%
 \at@end@of@kframe}
\makeatother

\definecolor{shadecolor}{rgb}{.97, .97, .97}
\definecolor{messagecolor}{rgb}{0, 0, 0}
\definecolor{warningcolor}{rgb}{1, 0, 1}
\definecolor{errorcolor}{rgb}{1, 0, 0}
\newenvironment{knitrout}{}{} % an empty environment to be redefined in TeX

\usepackage{alltt}
\usepackage{hyperref}
\signature{Ignacio Alvarez-Castro}
\address{}
\date{Montevideo, 4 de Marzo de 2024}
\IfFileExists{upquote.sty}{\usepackage{upquote}}{}
\begin{document}

\begin{letter}{Agencia Nacional de Investigación e Innovación}
\opening{Comisión Honoraria:}

A través de este medio, solicito extender el plazo para ejecutar los fondos del proyecto FSE\_S\_2022\_1\_173061 perteneciente al Fondo Sectorial de Energía 2022,
\textbf{por 6 meses, hasta el 31 de Agosto de 2024}. 

Los motivos para esta solicitud consisten en el atraso de actividades del proyecto debido a demoras en el acceso a datos y dificultades para la contratación de investigadores. En primer lugar, luego de reuiones con referentes del proyecto se obtuvo una parte de la información a fines del mes de Octubre/2023, siendo la única instancia en que recibimos información. Por otra parte, no pudimos contratar ayudantes con la formación necesaria y disponibilidad horaria para trabajar en el proyecto durante los meses de Abril a Noviembre de 2023, y no hubiera sido muy conveniente debido a que no teníamos los datos aun.

Desde Octubre el intercambio con referentes ha sido más fluido, permitiendo avanzar los primeros pasos en el análisis de la información. En Diciembre/2023 comenzó a trabajar en el proyecto un ayudante, estudiante avanzado de la licenciatura en Estadística, con cargo a los fondos del proyecto, y en Enero/2024 han comenzado a trabajar dos pasantes financiados por la universidad de Lyon.  Lo anterior muestra como las actividades del proyecto sufrieron un atraso importante y recién hemos comenzado a trabajar en los últimos meses de 2023. Desde ese momento, se mejoró el intercambio con al menos un referente de UTE, se conformó el equipo de trabajo y se ha comenzado a trabajar en 2 lineas que fueron identificadas como prioritarias para UTE. 

Por estos motivos se solicita extender el plazo del proyecto al 31 de Agosto de 2024, y en esa fecha evaluar el nivel de avance de los objetivos del proyecto.  


\closing{Sin otro particular,}
\begin{center}
\includegraphics[scale=.5]{../firma_nacho}
\end{center}

\end{letter}
\end{document}
